% !TeX program = xelatex
\documentclass[a4paper,10pt, twocolumn]{article}
\usepackage[a4paper,top=2cm,bottom=3cm,left=2cm,right=2cm,footskip=0.8cm]{geometry}
\usepackage{iftex}
\ifPDFTeX
  \usepackage[T1]{fontenc}
  \DeclareUnicodeCharacter{2212}{\ensuremath{-}}
  \usepackage{mathptmx,helvet}
\else
  \usepackage{fontspec}
  \IfFontExistsTF{Times New Roman}{\setmainfont{Times New Roman}}{\setmainfont{TeX Gyre Termes}}
  \IfFontExistsTF{Arial}{\setsansfont{Arial}}{\setsansfont{TeX Gyre Heros}}
\fi
\usepackage[brazilian]{babel}
\usepackage{csquotes,microtype,tikz,booktabs}
\usepackage[hidelinks]{hyperref}
\usepackage{graphicx}
\usepackage{float}
\usepackage{caption}    % Para personalizar legendas, opcional
\usepackage{xcolor}
\usepackage{fancyhdr}
\usepackage{titlesec}     % Para personalizar \section
\usepackage[style=ieee, sorting=none,maxbibnames=3,minbibnames=1, language=portuguese]{biblatex}
\definecolor{grayFEI}{HTML}{7F7F7F}
\usepackage{indentfirst}  % Indenta o primeiro parágrafo
\setlength{\parindent}{0.5cm}  % Aumenta a indentação para 2em (pode ser ajustado
\titlespacing*{\section}{0pt}{*0}{0pt}
\usepackage{ragged2e}


%-------------------------- ESTILOS NÃO ALTERAR ------------------------------------------------%
% Estilo para legendas de figura e tabela
\captionsetup[figure]{
  labelfont={color=azulFEI, bf, small, sf}, % "Figura" azul, negrito, Arial-like
  textfont={small, sf},                    % Texto da legenda em Arial-like
  font=footnotesize,                       % Tamanho aproximado de 9pt
  labelsep=period,                         % "Figura 1." e não "Figura 1:"
  name=Figura,                             % Tradução para "Figura"
  justification=raggedright,              % Alinhado à esquerda
  singlelinecheck=false                   % Garante alinhamento mesmo em uma linha
}

\captionsetup[table]{
  labelfont={color=azulFEI, bf, small, sf},
  textfont={small, sf},
  font=footnotesize,
  labelsep=period,
  name=Tabela,
  justification=raggedright,
  singlelinecheck=false
}

% Formatação section e subsection
\titleformat{\subsection}
  {\color{grayFEI}\normalfont\sffamily\selectfont\itshape\fontsize{10}{12}\selectfont}
  {\thesubsection.}{0.2em}{}
\renewcommand{\thesubsection}{\Alph{subsection}}

\titleformat{\section}
  {\color{azulFEI}\normalfont\sffamily\selectfont\bfseries\fontsize{11}{13}\selectfont} % Fonte Helvetica, tamanho 11pt, negrito
  {\thesection.}{0.2em}{}
\renewcommand{\thesection}{\Roman{section}} 
\titlespacing*{\section}
{0pt}{10pt}{0pt}

%Formatação Título
\newcommand{\inovatitle}[1]{%
  \noindent{\sffamily\bfseries\fontsize{18}{20}\selectfont\textcolor{azulFEI}{#1}\par}\vspace{12pt}
}

%Formatação Autores
\newcommand{\autores}[1]{
  \helvet{\fontsize{11}{13}\selectfont\textbf{\textit{#1}}}
}
\newcommand{\orientador}[1]{
  \helvet{\fontsize{11}{13}\selectfont\textbf{\textit{#1}}}
}
\newcommand{\filiacao}[1]{
  \helvet{\fontsize{10}{12}\selectfont{\textit{#1}}}
}
\newcommand{\email}[1]{
  \helvet{\fontsize{10}{12}\selectfont{\textit{#1}}}
}
\newcommand{\resumo}[1]{%
\vspace{15pt}
  \noindent
  {\helvet{\fontsize{12}{20}\selectfont\textbf{\textcolor{azulFEI}{Resumo:}}}} %
  \justifying #1\par\smallskip %
  \par
}

\newcommand{\keyword}[1]{
  \noindent\helvet{\fontsize{12}{20}\selectfont\textbf{\textcolor{azulFEI}{Palavras-chaves:}}} #1  \par\vspace{12pt}
}

\renewcommand{\thetable}{\Roman{table}}


% Define a cor azul específica (exemplo: azul FEI)
\definecolor{azulFEI}{HTML}{0070C0} % substitua pelo código que quiser

%Formatação Cabeçalho e rodapé
\setlength{\headsep}{1cm}
\newcommand{\helvet}[1]{{\sffamily\selectfont #1}}
\pagestyle{fancy}
\fancyhead{} % limpa cabeçalho
\renewcommand{\headrulewidth}{0.3pt} % ← ATIVA a linha (espessura ajustável)
\setlength\headheight{1.2cm}
\lhead{\includegraphics[width=4.97cm, height=1cm]{logoinova.png}}
\newcommand{\semester}{2º semestre/2026}
\rhead{\helvet{\semester}}

\fancyfoot{} % limpa rodapé
\lfoot{
  \sffamily\fontsize{8}{10}\selectfont
  Anais dos Trabalhos de Conclusão de Curso do Centro Universitário FEI, São Bernardo do Campo – SP\\
  \vspace{-4pt}\rule{\textwidth}{1pt} % ← linha abaixo do texto do rodapé
}
\rfoot{\sffamily\fontsize{8}{10}\selectfont \thepage}

%-------------------------------- FIM ESTILOS ---------------------------------------------------%



%Adicionar a bibliografia
\addbibresource{referencias.bib}


\begin{document}
\setlength{\columnsep}{20pt}
\raggedbottom
\twocolumn[
\begin{@twocolumnfalse}
\raggedright
\inovatitle{Um \textit{Harness} de IA para Resolução de Perguntas em Linguagem Natural: Adaptando o \textit{Ralph Wiggum Loop} Além do Desenvolvimento de \textit{Software}}

\autores{Hugo Emílio Nomura, Pedro Henrique Correia de Oliveira,\\
Pedro Henrique Satoru Lima Takahashi, Vitor Monteiro Vianna}\par
\filiacao{Ciência da Computação -- Centro Universitário FEI, \textit{campus} São Paulo}\par
\email{unifhnomura@fei.edu.br, unifpoliveira@fei.edu.br, unifptakahashi@fei.edu.br, unievvianna@fei.edu.br}\par\smallskip
\orientador{Orientador: Charles Henrique Porto Ferreira}\par
\filiacao{Departamento de Ciência da Computação, Centro Universitário FEI}\par
\email{cferreira@fei.edu.br}\par

\resumo{Adaptamos o \textit{Ralph Wiggum Loop}, originalmente concebido para o desenvolvimento de \textit{software}, para responder a perguntas em linguagem natural. O modelo divide a resolução em tarefas e verifica os resultados antes de formular a resposta. Nas questões científicas de pós-graduação do GPQA Diamond e no vestibular do ITA, a adaptação superou a inferência direta em 6,51 e 10,26 pontos percentuais (p.p.), respectivamente, quando ambas produziram respostas válidas. Contando as falhas como erros, o ganho persistiu no ITA (+4,76 p.p.), enquanto a acurácia no total de 743 questões ficou 2,02 p.p. abaixo da inferência direta.}
\keyword{\textit{Large Language Models} (LLMs). \textit{Ralph Wiggum Loop}. \textit{Harness} de IA. Engenharia de \textit{Prompt}. Decaimento de Contexto (\textit{context rot}).}
\end{@twocolumnfalse}
]

\section{Introdução}

No \textit{Ralph Wiggum Loop} original, o agente recebe uma solicitação de desenvolvimento de \textit{software}, como \enquote{adicione a recuperação de senha ao sistema}. Ele divide a implementação em tarefas com critérios de aceite e testes definidos. A cada tarefa, abre uma nova janela de contexto, lê os arquivos do projeto, altera o código, executa os testes e registra o que aprendeu~\cite{huntley}.

Como perguntas em linguagem natural não vêm com testes executáveis, o \textit{harness} usa o próprio modelo para definir critérios de aceite, verificar resultados parciais e compor a resposta. Comparamos sua acurácia à de uma chamada direta ao mesmo modelo.

\section{Metodologia}

No Planejamento, o modelo define os critérios gerais para a resposta final e divide a pergunta em tarefas. Cada enunciado traz os dados necessários à tarefa e seus critérios de aceite. O validador verifica se o plano é fiel ao enunciado, evita antecipar resultados, apresenta critérios verificáveis e cobre toda a pergunta sem redundâncias. Se houver problemas, o modelo refaz o plano com base nos ajustes apontados até obter aprovação ou esgotar as tentativas.

Na Execução em Loop, cada nova tarefa ou etapa começa em uma janela de contexto vazia, sem nenhum contexto anterior. Na primeira tentativa, só os aprendizados da questão, como unidades de medida, observações e resultados aprovados, complementam a descrição da tarefa. Isso busca mitigar o \textit{context rot}, a perda de qualidade associada à dispersão de informações em conversas longas.

O modelo resolve a tarefa e justifica a solução. Outra chamada verifica a validade técnica e os critérios de aceite. Em caso de reprovação, o modelo recebe as tentativas anteriores e os motivos da rejeição e refaz a tarefa.

\newpage
A quinta reprovação da mesma tarefa encerra a execução com erro. A temperatura inicial é zero para reduzir a variação das respostas e sobe para 1 ao refazer uma tarefa. O plano também admite cinco tentativas, com temperatura 0,5 na primeira revisão e 1 nas seguintes. Esse aumento introduz aleatoriedade e ajuda o modelo a explorar outra solução em vez de repetir o erro.

Na Criação da Resposta Final, o modelo verifica se os resultados aprovados atendem aos critérios gerais do Planejamento. Se houver lacunas, solicita tarefas ou revisões. Após a aprovação, redige a resposta e confere a alternativa.

\section{Metodologia de avaliação}

Usamos o modelo DeepSeek-V4-Flash-0731~\cite{deepseek} por seus pesos abertos, desempenho satisfatório, oferta por vários provedores e baixo custo por chamada. Para cada questão, comparamos duas execuções com enunciado, alternativas e instruções idênticos. Na direta, o modelo respondeu em uma única chamada e pôde justificar a escolha. Na orquestrada, o mesmo modelo executou as três fases do \textit{harness}. O modo de raciocínio nativo (\textit{thinking}) ficou desativado nos dois fluxos. Só a alternativa final foi comparada ao gabarito.

Das 753 questões previstas, 743 tinham resultados dos dois fluxos. Foram 268 do ENEM, da FUVEST, do ITA e da UNICAMP, 285 do MMLU~\cite{mmlu} (cinco por disciplina) e 190 do GPQA Diamond~\cite{gpqa}, com questões de física, química e biologia de pós-graduação. Nos dois últimos conjuntos, o modelo não recebeu exemplos de resolução.

Na avaliação completa, falhas operacionais contam como erros. Nos pares sem falha, ambos os fluxos entregaram uma alternativa válida, mesmo que incorreta. Comparamos os acertos pelo teste exato bilateral de McNemar e somamos os \textit{tokens} de entrada e saída de cada fluxo.

\clearpage

\section{Resultados e discussão}

\subsection{Acurácia por conjunto e por área}
Nos 743 pares, o \textit{harness} corrigiu 36 erros da inferência direta, mas errou 51 questões que ela havia acertado. Nos 711 pares sem falha, houve 35 correções e 26 perdas. Nenhuma das diferenças globais atingiu significância a 5\% (Tabela~\ref{tab:global}).

\begin{table}[H]
\caption{Acertos por configuração de inferência e valores de $p$ do teste exato bilateral de McNemar.}\label{tab:global}
\centering
\small
\begin{tabular}{@{}cccc@{}}
\toprule
Recorte & Direta & Orquestrada & $p$ \\
\midrule
Todos os pares & 652/743 & 637/743 & 0,1329 \\
Pares sem falha & 627/711 & 636/711 & 0,3057 \\
\bottomrule
\end{tabular}
\end{table}

No GPQA Diamond, a exclusão das falhas inverteu o saldo de acertos (Figura~\ref{fig:conjuntos}). Nos 21 pares excluídos por falha, a inferência direta acertou 17 questões e o \textit{harness} não acertou nenhuma. Nos 169 restantes, a orquestração acertou 134 e a direta, 123. Entre os pares sem falha dos demais conjuntos, o saldo foi de quatro acertos a mais nos vestibulares e seis a menos no MMLU.

\begin{figure}[H]
\centering
% Adaptação da Figura 1 de analise/GRAFICOS_RESULTADOS_PAPER.md para uma coluna.
\begin{tikzpicture}[x=0.40cm,y=0.65cm,font=\sffamily\footnotesize]
  \definecolor{completa}{HTML}{545F6A}
  \definecolor{condicional}{HTML}{0070A8}
  \draw[completa,fill=completa!23,line width=.4pt] (-8.2,3.8) rectangle ++(.5,.24);
  \node[anchor=west] at (-7.65,3.92) {Todos os pares};
  \fill[condicional] (0.1,3.8) rectangle ++(.5,.24);
  \node[anchor=west] at (.65,3.92) {Pares sem falha};
  \draw[black!55,thin] (0,-.50) -- (0,3.36);
  % Nome, posição vertical, diferenças de acurácia e rótulos decimais.
  \foreach \grupo/\y/\total/\valido/\a/\b in
    {Global/3/-2.02/1.27/{−2,02}/{+1,27},
     GPQA/2/-3.16/6.51/{−3,16}/{+6,51},
     MMLU/1/-2.81/-2.13/{−2,81}/{−2,13},
     Vestibulares/0/-.37/1.54/{−0,37}/{+1,54}} {
    \node[anchor=east] at (-5.45,\y) {\grupo};
    \draw[completa,fill=completa!23,line width=.4pt] (0,\y+.10) rectangle (\total,\y+.31);
    \fill[condicional] (0,\y-.31) rectangle (\valido,\y-.10);
    \node[anchor=east,inner sep=1pt] at (\total,\y+.20) {\a};
    \ifdim\valido pt<0pt
      \node[anchor=east,inner sep=1pt] at (\valido,\y-.20) {\b};
    \else
      \node[anchor=west,inner sep=1pt] at (\valido,\y-.20) {\b};
    \fi
  }
\end{tikzpicture}
\caption{Diferença de acurácia em pontos percentuais. Valores positivos indicam maior acurácia da inferência orquestrada. Houve respostas válidas nas duas configurações em 169 dos 190 pares do GPQA, 282 dos 285 do MMLU e 260 dos 268 dos vestibulares.}\label{fig:conjuntos}
\end{figure}

Entre os pares sem falha, o saldo foi de 14 acertos a mais em exatas e tecnologia (372 questões), dois a mais em biológicas e saúde (101), sete a menos em humanas e sociais (228) e zero nas interdisciplinares (10). Em exatas e tecnologia, o valor de $p$ ajustado por Holm para as quatro comparações foi 0,134.

Ao agrupar os pares sem falha pelo consumo de \textit{tokens} da inferência direta dentro de cada conjunto, observamos que a diferença de acurácia passou de $-0{,}55$ ponto percentual no quartil de menor consumo para $+3{,}39$ no de maior consumo.

\newpage
\subsection{Confiabilidade da validação}
O \textit{harness} registrou 31 falhas operacionais e a inferência direta, duas, uma delas em comum. Em 16 das 31 falhas, a mesma tarefa foi reprovada cinco vezes. Oito decorreram de saídas em formato inválido. Também houve erros de conteúdo. Em 20 dos 26 pares sem falha em que só a direta acertou, o verificador aprovou todas as tarefas na primeira tentativa.

Uma questão do ITA de 2024 pedia um intervalo ao qual o raio de um cilindro \emph{poderia} pertencer. O modelo deduziu $r>2$ e escolheu corretamente \mbox{$2<r<2{,}5$}. O plano, porém, exigia que a alternativa abrangesse \emph{todos} os raios possíveis. Após cinco reprovações, a execução terminou sem resposta. O caso mostra como um critério mal formulado no planejamento pode levar o verificador a rejeitar até a solução correta.

\subsection{Custo e alcance dos resultados}
Nos 711 pares sem falha, a inferência orquestrada consumiu 41,80 milhões de \textit{tokens}, ou 98,36 vezes o consumo da direta (425 mil). Cada um dos nove acertos líquidos custou, em média, 4,60 milhões de \textit{tokens} adicionais. A duração mediana por questão passou de 1,75 para 37,24 segundos.

Avaliou-se um modelo, uma vez por fluxo e questão, sem igualar o orçamento de \textit{tokens} nem isolar o efeito de etapas como a renovação do contexto.

\section{Conclusão}

Nos 711 pares em que ambos os fluxos produziram respostas válidas, o \textit{Ralph Wiggum Loop} adaptado obteve nove acertos a mais que a inferência direta (+1,27 p.p.). Ao incluir as falhas operacionais, o saldo se inverteu. Nas 743 questões, foram 15 acertos a menos ($-2{,}02$ p.p.).

No GPQA Diamond, com questões científicas de pós-graduação em nível de doutorado~\cite{gpqa}, o ganho foi de 6,51 p.p. nos 169 pares sem falha. O conjunto concentrou 21 das 31 falhas do \textit{harness}, com perda de 3,16 p.p. no total. No vestibular do ITA, a vantagem alcançou 10,26 p.p. nos 39 pares sem falha e persistiu nas 42 questões (+4,76 p.p.). A comparação é exploratória, pois não há uma escala comum de dificuldade. Ainda assim, os ganhos nesses dois testes sugerem que a adaptação favorece a resolução de questões de maior complexidade técnica quando o fluxo produz uma resposta válida.

\section{Referências}
\printbibliography[heading=none]
\end{document}
